(1) 令 \(J=I_0+m\ell^2\) 为粘连后的转动惯量。极短碰撞期间轴承冲量作用于转轴，对该轴的力矩冲量为零，重力矩冲量忽略，因此有 \(mv\ell=J\Omega_0\)，即 \(\Omega_0=mv\ell/J\)；线动量不因固定轴反力而守恒。
(2) 碰撞前动能 \(\frac12mv^2\)，碰撞后 \(\frac12J\Omega_0^2=\frac{m^2v^2\ell^2}{2J}\)，故损失 \(\Delta E=\frac12mv^2(1-m\ell^2/J)\ge0\)。设 \(G=Mg d+mg\ell\)，碰撞后的无耗散上升满足 \(\frac12J\Omega_0^2=G(1-\cos\theta_{\max})\)，当左侧小于 \(2G\) 时有唯一 \(0\le\theta_{\max}<\pi\)。
(3) 粘连后的引力势能为 \(G(1-\cos\theta)\)，动能为 \(J\dot\theta^2/2\)；Euler–Lagrange 方程为 \(J\ddot\theta+G\sin\theta=0\)。在 \(\theta=0\) 附近线性化给 \(\omega_{\rm small}=\sqrt{G/J}\)，当 \(G,J>0\) 时稳定，小角振幅虽任意小但不能由此判断大角周转。
(4) 刚性铰链允许推拉，越过顶部仅要求碰撞后能量严格大于顶部位能 \(2G\)，即 \(\Omega_0^2>4G/J\)；等号恰停在不稳定平衡的分界线上。碰撞时守恒的是关于固定轴的角动量，粘连必有内耗，绝不能把弹丸碰前动能等同于碰后转动动能。
